% ============================================================
% CAS/Python ενότητες — Κεφάλαιο 2: Άλγεβρα Πινάκων
%
% QR code: qr_cb01e79b72b4.png  (→ latex/qrcodes/)
% URL: https://nikmatz.github.io/ELADIC_GR/code/ch02/
% ============================================================

\casactivbox%
  {Maxima: Πράξεις Πινάκων — Γινόμενο, Αντίστροφος, LU}%
  {%
    \label{cas:ch2_matrix_algebra}
    Δίνονται οι πίνακες
    \[
      A = \begin{pmatrix} 1 & 2 & 3 \\ 0 & 1 & 4 \\ 5 & 6 & 0 \end{pmatrix},
      \quad
      B = \begin{pmatrix} 2 & -1 & 0 \\ 1 & 3 & 1 \\ 0 & 2 & 4 \end{pmatrix},
      \quad
      \mathbf{b} = \begin{pmatrix} 1 \\ 2 \\ 3 \end{pmatrix}.
    \]
    \begin{enumerate}[label=(\alph*),itemsep=2pt]
      \item Υπολογίστε $A+B$, $3A$, $AB$ και $BA$. Επαληθεύστε ότι
        $AB \neq BA$.
      \item Υπολογίστε $\det(A)$, $\det(B)$ και επαληθεύστε
        $\det(AB) = \det(A)\cdot\det(B)$.
      \item Βρείτε $A^{-1}$ και επαληθεύστε $A A^{-1} = I_3$.
      \item Λύστε το σύστημα $A\mathbf{x} = \mathbf{b}$ χρησιμοποιώντας
        \texttt{invert(A)} και επαληθεύστε.
    \end{enumerate}
    \smallskip
    \textbf{Εντολές Maxima:}
    \texttt{transpose(A)}, \texttt{determinant(A)},
    \texttt{invert(A)}, \texttt{A . B}, \texttt{rref(A)}
  }%
  {https://nikmatz.github.io/ELADIC_GR/code/ch02/}%
  {qr_cb01e79b72b4.png}

\subsection*{Δραστηριότητα Python}

\begin{pybox}{Python: Πίνακες — NumPy, Αποσύνθεση LU, Heatmaps}%
             {https://nikmatz.github.io/ELADIC_GR/code/ch02/}%
             {qr_cb01e79b72b4.png}
\label{py:ch2_matrix_algebra}
Χρησιμοποιήστε τους πίνακες $A$, $B$ και το $\mathbf{b}$ της δραστηριότητας
Maxima.
\begin{enumerate}[label=(\alph*),itemsep=2pt]
  \item Εκτελέστε τις πράξεις με \texttt{numpy}: \texttt{A @ B},
    \texttt{A.T}, \texttt{np.linalg.det(A)}, \texttt{np.linalg.inv(A)}.
    Συγκρίνετε με τα αποτελέσματα Maxima.
  \item Λύστε $A\mathbf{x}=\mathbf{b}$ με \texttt{np.linalg.solve}
    και επαληθεύστε αριθμητικά.
  \item Εκτελέστε την αποσύνθεση $A = PLU$ με \texttt{scipy.linalg.lu}.
    Επαληθεύστε $P \cdot L \cdot U = A$.
  \item Για τον συμμετρικό πίνακα
    $S = \begin{pmatrix}4&2&1\\2&5&3\\1&3&6\end{pmatrix}$,
    βρείτε τις ιδιοτιμές με \texttt{np.linalg.eigh}
    και επαληθεύστε ότι $S$ είναι θετικά ορισμένος.
  \item \textit{Προαιρετικό:} Οπτικοποιήστε $A$, $AB$ και $L$ ως
    heatmaps (Matplotlib \texttt{imshow}).
\end{enumerate}
\medskip
\textbf{Βιβλιοθήκες / Εντολές:}
\texttt{A @ B}, \texttt{np.linalg.det()}, \texttt{np.linalg.inv()},
\texttt{np.linalg.solve()}, \texttt{scipy.linalg.lu()},
\texttt{np.linalg.eigh()}
\end{pybox}
